% Options for packages loaded elsewhere
\PassOptionsToPackage{unicode}{hyperref}
\PassOptionsToPackage{hyphens}{url}
\documentclass[
  ignorenonframetext,
]{beamer}
\newif\ifbibliography
\usepackage{pgfpages}
\setbeamertemplate{caption}[numbered]
\setbeamertemplate{caption label separator}{: }
\setbeamercolor{caption name}{fg=normal text.fg}
\beamertemplatenavigationsymbolsempty
% remove section numbering
\setbeamertemplate{part page}{
  \centering
  \begin{beamercolorbox}[sep=16pt,center]{part title}
    \usebeamerfont{part title}\insertpart\par
  \end{beamercolorbox}
}
\setbeamertemplate{section page}{
  \centering
  \begin{beamercolorbox}[sep=12pt,center]{section title}
    \usebeamerfont{section title}\insertsection\par
  \end{beamercolorbox}
}
\setbeamertemplate{subsection page}{
  \centering
  \begin{beamercolorbox}[sep=8pt,center]{subsection title}
    \usebeamerfont{subsection title}\insertsubsection\par
  \end{beamercolorbox}
}
% Prevent slide breaks in the middle of a paragraph
\widowpenalties 1 10000
\raggedbottom
\AtBeginPart{
  \frame{\partpage}
}
\AtBeginSection{
  \ifbibliography
  \else
    \frame{\sectionpage}
  \fi
}
\AtBeginSubsection{
  \frame{\subsectionpage}
}
\usepackage{iftex}
\ifPDFTeX
  \usepackage[T1]{fontenc}
  \usepackage[utf8]{inputenc}
  \usepackage{textcomp} % provide euro and other symbols
\else % if luatex or xetex
  \usepackage{unicode-math} % this also loads fontspec
  \defaultfontfeatures{Scale=MatchLowercase}
  \defaultfontfeatures[\rmfamily]{Ligatures=TeX,Scale=1}
\fi
\usepackage{lmodern}
\ifPDFTeX\else
  % xetex/luatex font selection
\fi
% Use upquote if available, for straight quotes in verbatim environments
\IfFileExists{upquote.sty}{\usepackage{upquote}}{}
\IfFileExists{microtype.sty}{% use microtype if available
  \usepackage[]{microtype}
  \UseMicrotypeSet[protrusion]{basicmath} % disable protrusion for tt fonts
}{}
\makeatletter
\@ifundefined{KOMAClassName}{% if non-KOMA class
  \IfFileExists{parskip.sty}{%
    \usepackage{parskip}
  }{% else
    \setlength{\parindent}{0pt}
    \setlength{\parskip}{6pt plus 2pt minus 1pt}}
}{% if KOMA class
  \KOMAoptions{parskip=half}}
\makeatother
\setlength{\emergencystretch}{3em} % prevent overfull lines
\providecommand{\tightlist}{%
  \setlength{\itemsep}{0pt}\setlength{\parskip}{0pt}}
\usepackage{bookmark}
\IfFileExists{xurl.sty}{\usepackage{xurl}}{} % add URL line breaks if available
\urlstyle{same}
\hypersetup{
  hidelinks,
  pdfcreator={LaTeX via pandoc}}

\author{\texorpdfstring{}{}}
\date{}

\begin{document}

\begin{frame}[fragile]{API and Workflows}
\protect\phantomsection\label{api-and-workflows}
This section describes the programmatic surface that the shipped Python
package exposes in practice, aligned with
\texttt{../cogant/docs/API\_GUIDE.md} and
\texttt{../cogant/VERIFICATION\_REPORT.md}. It is \textbf{not} an
empirical benchmark section; COGANT does not ship comparative timing
claims in the manuscript layer. Instead, it records the \textbf{surface
area} users can rely on: two complementary entry points (a Session for
stepwise work and a Pipeline for batch runs), the Bundle accessors that
expose their artifacts, a command-line interface, and a Review API for
human-in-the-loop curation.

\begin{block}{Session-oriented workflow}
\protect\phantomsection\label{session-oriented-workflow}
The Session API is intended for interactive exploration: it is the right
choice when a user wants to inspect intermediate artifacts between
stages, iterate on a single repository in a notebook, or debug a
specific extraction step without committing to a full scripted run. Each
call returns control to the caller so that the graph, mappings, or
state-space model can be examined before the next stage is invoked.

\texttt{Session.from\_target} accepts a local path or URL, then supports
a stepwise workflow:

\begin{itemize}
\tightlist
\item
  \texttt{extract\_static} --- AST-oriented extraction for supported
  languages.
\item
  \texttt{extract\_dynamic} --- traces and coverage when inputs exist.
\item
  \texttt{build\_graph} --- program graph construction.
\item
  \texttt{translate\_to\_gnn} --- Generalized Notation Notation (GNN)
  representation.
\item
  \texttt{compile\_state\_space} --- behavioral model when the pipeline
  has sufficient data.
\item
  \texttt{export\_all} --- writes JSON artifacts under a chosen output
  directory.
\end{itemize}

This path suits interactive notebooks and incremental debugging.
\end{block}

\begin{block}{Pipeline-oriented workflow}
\protect\phantomsection\label{pipeline-oriented-workflow}
The Pipeline API is intended for scripted, reproducible batch runs: it
is the right choice when a user wants to process many repositories with
a fixed configuration, wire COGANT into CI, or guarantee that every run
executes the same ordered stages with the same plugin settings. Because
the whole run is described by a single \texttt{PipelineConfig}, it can
be checked into version control and replayed without manual
intervention.

\texttt{PipelineRunner} with \texttt{PipelineConfig} runs an ordered
list of stages (ingest, static, normalize, graph, translate, state
space, process, export, validate). Configuration can skip stages (for
example dynamic analysis), attach \textbf{plugin} settings per language,
set \texttt{output\_dir}, verbosity, and dry-run mode. Results aggregate
into a \textbf{Bundle} with \texttt{stage\_results}, error lists, and
accessors described below.
\end{block}

\begin{block}{Bundle accessors}
\protect\phantomsection\label{bundle-accessors}
The bundle API exposes summaries and artifacts, including:

\begin{itemize}
\tightlist
\item
  \texttt{repo\_summary}, \texttt{program\_graph},
  \texttt{state\_space\_model}, \texttt{process\_model}
\item
  \texttt{gnn\_markdown} --- human-readable graph summary
\item
  \texttt{validation\_report}
\item
  \texttt{render\_site} --- static HTML site with graph and model views
\item
  \texttt{to\_json} / \texttt{save\_json}
\end{itemize}

The verification report lists these as delivered capabilities for the
v0.1.0 line.
\end{block}

\begin{block}{Command-line interface}
\protect\phantomsection\label{command-line-interface}
The CLI entry point (\texttt{cogant.cli.main}) implements commands such
as \texttt{init}, \texttt{scan}, \texttt{extract-static},
\texttt{extract-dynamic}, graph and export verbs, and
validation-oriented subcommands. Exact flags live in
\texttt{../cogant/docs/CLI\_GUIDE.md}; the manuscript does not duplicate
them to avoid drift.
\end{block}

\begin{block}{Review API}
\protect\phantomsection\label{review-api}
\texttt{ReviewAPI} supports interactive curation: load a bundle, present
mappings, accept, reject, or edit, then save a curated bundle. This
closes the loop when human review is part of the ML dataset
construction.
\end{block}
\end{frame}

\end{document}
